\documentclass{article}

\usepackage[preprint]{neurips_2025}

\usepackage[utf8]{inputenc}
\usepackage[T1]{fontenc}
\usepackage{hyperref}
\usepackage{url}
\usepackage{booktabs}
\usepackage{amsfonts}
\usepackage{amsmath}
\usepackage{nicefrac}
\usepackage{microtype}
\usepackage{xcolor}
\usepackage{enumitem}
\usepackage{graphicx}
\usepackage{subcaption}
\usepackage{algorithm}
\usepackage{algorithmic}

\title{RL-based Dynamic Scheduling for LLM Serving:\\Midterm Progress Report}

\author{
  Yuxuan Bai \\
  University of Virginia \\
  \texttt{yuxuan.bai@virginia.edu} \\
  \And
  Jiaqi Liu \\
  University of Virginia \\
  \texttt{jiaqi.liu@virginia.edu} \\
  \And
  Peng Xia \\
  University of Virginia \\
  \texttt{peng.xia@virginia.edu} \\
}

\begin{document}

\maketitle

\begin{abstract}
% TODO: Write abstract summarizing midterm progress
% 1 paragraph: problem, approach, progress so far, preliminary findings
\end{abstract}

\section{Introduction}

% TODO: Problem statement (0.5 page)
% - LLM serving scheduling challenges
% - Llumnix and its limitations (heuristic-based)
% - Our approach: RL-based virtual usage replacement
% - Key design: minimal modification to Llumnix

\section{Background}

% TODO: Background (1 page)

\subsection{LLM Serving and KV Cache Management}
% Auto-regressive inference, KV cache growth, preemption problem

\subsection{Llumnix System}
% Virtual usage, freeness, 4 heuristic rules, migration mechanism

\subsection{RL for Systems Scheduling}
% Prior work: Pensieve, Decima, resource management with RL

\section{Progress}

% TODO: Progress so far (1.5-2 pages)

\subsection{MDP Formulation}

% State space definition with math
% Action space
% Reward function

\paragraph{State space.}
% Per-instance: memory usage, queue length, waiting demand, etc.
% Cluster: mean load, variance, total waiting, etc.

\paragraph{Action space.}
% Continuous virtual usage adjustment per instance

\paragraph{Reward function.}
% r_t = w_1 * latency + w_2 * preemption + w_3 * fragmentation + ...

\subsection{Simulation Environment}
% Description of the simulator
% Timing model based on Llumnix paper data
% How it interfaces with Llumnix scheduling logic

\subsection{RL Agent Architecture}
% Policy network design
% Multi-head attention architecture
% PPO training setup

% TODO: Include figure of network architecture
% \begin{figure}[h]
%   \centering
%   \includegraphics[width=0.8\textwidth]{figures/architecture.pdf}
%   \caption{RL policy network architecture for Llumnix scheduling.}
%   \label{fig:architecture}
% \end{figure}

\subsection{Preliminary Results}
% Training curves
% Initial comparison in simulation

% TODO: Include training curve figure
% \begin{figure}[h]
%   \centering
%   \includegraphics[width=0.6\textwidth]{figures/training_curve.pdf}
%   \caption{Training reward over timesteps.}
%   \label{fig:training}
% \end{figure}

% TODO: Include preliminary results table
% \begin{table}[h]
%   \centering
%   \caption{Preliminary comparison in simulation (ShareGPT trace, 16 instances).}
%   \begin{tabular}{lcccc}
%     \toprule
%     Policy & Request P99 (s) & Prefill P99 (s) & Decode P99 (s) & Preemptions \\
%     \midrule
%     Round-Robin & & & & \\
%     Llumnix-Heuristic & & & & \\
%     Llumnix-RL (ours) & & & & \\
%     \bottomrule
%   \end{tabular}
% \end{table}

\section{Remaining Work}

% TODO: What's left (0.5 page)
% - Full Llumnix integration on real GPUs
% - Complete experiment suite (5 experiments)
% - Ablation studies
% - Final report and presentation

\section{Challenges}

% TODO: Technical challenges (0.5 page)
% - Simulation fidelity vs. speed tradeoff
% - Reward shaping (what weights work best)
% - Integration complexity (code structure differs from paper description)

\begin{thebibliography}{6}
\small

\bibitem[Sun et al.(2024)]{sun2024llumnix}
B.~Sun, Z.~Huang, H.~Zhao, W.~Xiao, X.~Zhang, Y.~Li, and W.~Lin.
Llumnix: Dynamic scheduling for large language model serving.
In \emph{OSDI}, 2024.

\bibitem[Kwon et al.(2023)]{kwon2023vllm}
W.~Kwon et al.
Efficient memory management for large language model serving with {PagedAttention}.
In \emph{SOSP}, 2023.

\bibitem[Schulman et al.(2017)]{schulman2017proximal}
J.~Schulman, F.~Wolski, P.~Dhariwal, A.~Radford, and O.~Klimov.
Proximal policy optimization algorithms.
\emph{arXiv:1707.06347}, 2017.

\bibitem[Mao et al.(2016)]{mao2016resource}
H.~Mao, M.~Alizadeh, I.~Menache, and S.~Kandula.
Resource management with deep reinforcement learning.
In \emph{HotNets}, 2016.

\bibitem[Mao et al.(2017)]{mao2017neural}
H.~Mao, R.~Netravali, and M.~Alizadeh.
Neural adaptive video streaming with {Pensieve}.
In \emph{SIGCOMM}, 2017.

\end{thebibliography}

\end{document}
